\section{Điều bản đồ chưa giải quyết}
\label{sec:ch07-dieu-chua-giai-quyet}

\index{độ trung thực}%
\index{ground truth}%
\index{chỉ số đánh giá}%
Khung hình thức và bản đồ của khảo sát cho một chỗ đứng để nhìn cả lĩnh vực, và
chính chỗ đứng ấy làm lộ ra ba vấn đề mà không mục nào trước đây giải quyết
được. Bài báo quan điểm liệt kê cả ba như thách thức chung của mọi nhánh
attribution~\autocite{p18unified}.

Vấn đề thứ nhất là chi phí tính toán, và nó là hệ quả trực tiếp của kỹ thuật
dùng chung. Phép nhiễu vấp phải \tn{lời nguyền số chiều}{curse of
dimensionality} ở cả ba nhánh, vì số tập con
cần thử tăng theo hàm mũ. Gradient rẻ hơn nhưng chỉ cho thông tin bậc nhất, nên
muốn tinh hơn thì phải cộng gradient qua nhiều bước hoặc tính Hessian.
\tn{Xấp xỉ tuyến tính}{linear approximation} cần đủ nhiều lượt đánh giá mô hình
để khớp được một mô hình tuyến tính đáng tin~\autocite{p18unified}. Mỗi cách khắc phục đều là một phép cắt bớt,
và mỗi phép cắt bớt là một tham số mới.

Vấn đề thứ hai theo ngay sau: kết quả không ổn định giữa các lần chạy. Nguồn
ngẫu nhiên có ở khắp nơi trong bộ kỹ thuật chung, từ phép lấy mẫu Monte Carlo
trong nhóm phép nhiễu đến tối ưu ngẫu nhiên trong nhóm xấp xỉ tuyến tính. Chồng
lên đó là các \tn{siêu tham số}{hyperparameter} không tầm thường: số mẫu dùng
để tính giá trị Shapley,
tốc độ học và số bước của phép khớp tuyến tính, và lượng giảm chấn cho nghịch
đảo Hessian. Riêng với kiểu nhiễu được chọn, đặt về không, thay bằng trung
bình hay lấy ngẫu nhiên, bài báo nhận xét rằng ba cách lẽ ra phải cho kết quả
gần nhau theo nguyên lý, ở cả ba nhánh, nhưng trên thực tế thường cho
attribution khác nhau rõ rệt~\autocite{p18unified}. Chương~\ref{ch:lime-nao-dang-tin} đọc kỹ hiện
tượng này ở LIME.

Vấn đề thứ ba là vấn đề của cuốn sách. Muốn nói phương pháp nào tốt hơn thì phải
đánh giá, mà đánh giá attribution thì không có ground truth để đối
chiếu~\autocite{p18unified}. Ba lối thoát đang được dùng, và bài báo ghi giá của
từng lối. Đánh giá phản thực bỏ đi những phần tử điểm cao rồi đo mô hình tụt bao
nhiêu; nó cho một con số cụ thể, nhưng sinh ra các phản thực thì đắt, và tương
tác giữa các phần tử có thể không lọt vào phép đo từng phần tử một. Đánh giá
theo tác vụ đo lợi ích thực tế của lời giải thích trong một việc cụ thể, đổi lại
kết quả không chắc mang sang được tác vụ khác. Đánh giá bằng con người được coi
là chuẩn vàng, nhưng chủ quan, tốn kém, và bài báo chỉ nhận nó là chuẩn vàng
trong một số trường hợp~\autocite{p18unified}.

Chi tiết đáng giữ nhất nằm trong lối thứ nhất. Bài báo gọi tên ba đại lượng của
đánh giá phản thực: fidelity đo tính đủ, bằng cách giữ lại những phần tử điểm
cao và bỏ những phần tử điểm thấp; inverse fidelity đo tính cần, bằng cách làm
ngược lại; và sparsity đo cần bao nhiêu phần tử mới đạt fidelity
cao~\autocite{p18unified}. Leave-one-out của
\tn{attribution dữ liệu}{data attribution} là một trường hợp riêng của inverse
fidelity. Hai trong ba đại lượng ấy, fidelity và inverse fidelity, là hai họ
chỉ số mà chương~\ref{ch:do-do-trung-thuc} sẽ đọc kỹ dưới tên riêng của từng
bài báo, và ở đây chúng đến từ một bài báo không hề phê phán chúng.

Phép hợp nhất vì thế có hệ quả ở cả hai chiều. Nếu ba nhánh
thật sự dùng chung một bộ kỹ
thuật và một bộ chỉ số đánh giá, thì một sai sót trong bộ chỉ số ấy không dừng
lại ở một nhánh. Bản thân các tác giả cũng ghi lại quan điểm đối lập rằng ba
nhánh nên tiếp tục phát triển độc lập vì mỗi nhánh có động lực riêng, và họ trả
lời bằng lập luận chứ không bằng bằng chứng thực nghiệm~\autocite{p18unified}.
Phần~IV bắt đầu từ đúng chỗ đó: chương~\ref{ch:do-do-trung-thuc} xem xét các họ
chỉ số dùng để đo độ trung thực, và chương~\ref{ch:chi-so-khong-do-do-trung-thuc}
hỏi các chỉ số ấy có đo được thứ chúng khai báo hay không.
